\subsection{The $Q_8$ Prototype}
  \label{subsec:q8-prototype}

  To illustrate the admissibility mechanism explicitly, we consider the
  Cayley graph of the quaternion group $Q_8$.
  This example provides the simplest analytically tractable setting in which the
  bounded-flux constraint induces a non-trivial hierarchy among spectral sectors.

  The quaternion group
  \[
    Q_8=\{\pm 1,\pm i,\pm j,\pm k\}
  \]
  is generated by the elements $\{\pm i,\pm j,\pm k\}$ with relations
  \[
    i^2=j^2=k^2=ijk=-1 .
  \]

  We consider the undirected Cayley graph constructed from the generating set
  \[
    S=\{\pm i,\pm j,\pm k\}.
  \]
  Each vertex corresponds to a group element and edges connect elements related
  by right multiplication by a generator in $S$.

  The relational Laplacian of this graph admits a spectral decomposition indexed
  by the irreducible representations of $Q_8$.
  Besides four one-dimensional representations, $Q_8$ possesses a
  two-dimensional irreducible representation corresponding to its spinorial
  sector.

  Let $\lambda_\rho$ denote the Laplacian eigenvalue associated with the
  representation sector $\rho$.
  Applying the admissibility envelope derived in
  Section~\ref{subsec:admissibility-bounded-flux},
  \[
    A_\rho^{\max}=\frac{c_\chi}{\sqrt{\lambda_\rho}},
  \]
  one obtains the maximal amplitude that each sector can sustain without
  violating the bounded-flux constraint.

  A direct computation of the Laplacian spectrum shows that the admissible
  amplitude of the two-dimensional sector exceeds that of the one-dimensional
  sectors by a factor
  \[
    \sqrt{\frac{4}{3}} .
  \]

  Thus, even though several representation sectors are kinematically present in
  the graph, the bounded-flux constraint dynamically favours the non-abelian
  sector.

  In physical terms, the admissibility envelope acts as a spectral filter that
  suppresses sectors associated with larger Laplacian eigenvalues.
  The lowest non-trivial sector therefore becomes the dominant carrier of
  admissible fluctuations.

  This result should not be interpreted as a special property of the quaternion
  group itself.
  Rather, $Q_8$ provides the minimal relational structure in which the mechanism
  can be computed exactly.

  The strengthening of this hierarchy in larger binary groups, and in particular
  in the binary icosahedral group $2I$, is analysed in the following section.
